\scanpage{23}
\reversemarginpar
\margin{\handwritten{SGA 4 $\overline{\text{\uline{VI}}}$ 2}}
\normalmarginpar

\uline{Exercice} \handwritten{2}. . Soi\inserted{en}t $C$ une \deleted{\uline{$U$}-}\inserted{petite} catégorie\handwritten{,}\inserted{$C'$ une enveloppe de Karoubi\inserted{de $C$ (IV 7.5\handwritten{b))}}} Montrer que les conditions suivantes sont équivalentes:

(i) La catégorie $\mathrm{Ind}(C)$ (I 8.2) est un \uline{$U$}-topos.

( bis) Le foncteur canonique $\mathrm{Ind}(C)\handwritten{\longrightarrow}\overset{\handwritten{\scriptscriptstyle\wedge}}{C}$ admet un adjoint à droite qui \deleted{commute aux} est exact à gauche.

(i ter) Il existe une topologie sur $C$ pour laquelle les faisceaux soient les préfaisceaux ind-représentables.

% The outlined block is cancelled by four oblique strokes; the supplied
% deletion macro records this state, including its earlier internal changes.
\deleted{(ii) Dans $C'$, les limites inductives finies sont représentables, \deleted{et munissant}\inserted{il existe sur} $C'$ \deleted{\unsure{d'une}}\inserted{une} topologie dont les familles $(X_i\quad X)_{i\ I}$ sont celles qui sont raffinées par une famille épimorphique finie}

(ii) \deleted{Dans $C'$} I\begin{tikzpicture}[baseline=(letter.base),x=1em,y=1ex]
\node[inner sep=0pt,anchor=base] (letter) at (0,0) {l};
\draw[line width=.35pt] (-.22,1.85)--(.22,1.85)--(.22,2.05);
\draw[line width=.35pt] (-.18,-.42)--(.11,-.42)--(.11,-.65);
\end{tikzpicture} existe une topologie sur $C'$ dont les familles couvrantes sont celles qui so\unsure{\deleted{j}n|\deleted{n}j}t raffinées par une famille épimorphique (ou encore, épimorphique effective universelle ..) \uline{finie}, dans $C'$ les limites \deleted{projectives}\inserted{inductives} finies sont représentables, et le foncteur canonique $C\handwritten{\longrightarrow}\overset{\handwritten{\scriptscriptstyle\sim}}{C}$ \deleted{\unsure{est}} y commute.

(iii) (Si dans $C$ les limites projectives finies sont représentab\deleted{\unsure{l|e}}les,\deleted{\unsure{â|à}} de sorte qu'on peut prendre $C'=C$.) Dans $C$ les limites inductives finies sont représentables, et $C$ satisfait la condition ST) de \reversemarginpar\margin{\handwritten{$\times$}}\normalmarginpar I 8.8 f).

(iv) Il existe un topos $E$ engendré par sa sous-catégorie pleine \reversemarginpar\margin{\handwritten{$\times$}}\normalmarginpar $E_{\mathrm{PF}}$ (\quad\quad), et une équivalence de catégories $C'\handwritten{\approxeq}E_{\mathrm{PF}}$.

De plus, \deleted{la topologie} si ces conditions sont satisfaites, la topologie postulée dans (i ter) est nécessairement celle décrite dans (ii), et le topos $E$ postulé dans (iv) est nécessairement équivalent à $\mathrm{Ind}(C)$. Ce topos est donc parfait si et seulement si \inserted{dans} $C'$ \deleted{est} les limites projectives finies sont représentables.
% archive sheet number: 22

